\documentclass[mode=notes, palette=ocean, fontsize=11pt]{modernclassnotes}

\course{LaTeX Package Documentation}
\coursecode{MCN v1.0}
\professor{Modern Class Notes Maintainers}
\institution{Open Source TeX Initiative}
\term{Release 2026}
\docsubtitle{User Manual \& Architecture Guide for Class Notes Package}
\title{The \texttt{modernclassnotes} Class \& Package}

\begin{document}

\maketitle

\tableofcontents
\newpage

\section{Introduction}

The \texttt{modernclassnotes} package and document class provides a modern, highly aesthetic, and open-source LaTeX framework designed specifically for university professors, lecturers, and educators. It enables effortless creation and distribution of lecture notes, student handouts, and problem worksheets.

\begin{takeaways}
Key highlights of \texttt{modernclassnotes}:
\begin{itemize}
    \item \textbf{Three Specialized Modes}: \texttt{notes} (full course manuals), \texttt{handout} (1--2 page lecture summaries), and \texttt{worksheet} (homeworks and solution keys).
    \item \textbf{Curated Color Palettes}: Switch between 7 themes (\texttt{ocean}, \texttt{nord}, \texttt{midnight}, \texttt{emerald}, \texttt{burgundy}, \texttt{amethyst}, \texttt{mono}) with a single class option.
    \item \textbf{Zero Mandatory Heavy Dependencies}: Native TeX/LaTeX box engine ensuring maximum portability across TeX Live, MiKTeX, and Overleaf.
    \item \textbf{Lecture Tracking}: Automated lecture header banners, running headers, and table of contents logging via \texttt{\textbackslash lecture}.
\end{itemize}
\end{takeaways}

\section{Class Options \& Configuration}

Load the class in your LaTeX document header:
\begin{codebox}[Loading the Class]
\begin{lstlisting}[language=TeX]
\documentclass[mode=notes, palette=ocean, fontsize=11pt, font=default]{modernclassnotes}
\end{lstlisting}
\end{codebox}

\subsection{Available Class Keys}
\begin{table}[h]
\centering
\begin{tabular}{lll}
\toprule
\textbf{Key} & \textbf{Values} & \textbf{Default} \\
\midrule
\texttt{mode} & \texttt{notes}, \texttt{handout}, \texttt{worksheet} & \texttt{notes} \\
\texttt{palette} & \texttt{ocean}, \texttt{midnight}, \texttt{nord}, \texttt{emerald}, \texttt{burgundy}, \texttt{amethyst}, \texttt{mono} & \texttt{ocean} \\
\texttt{fontsize} & \texttt{10pt}, \texttt{11pt}, \texttt{12pt} & \texttt{11pt} \\
\texttt{font} & \texttt{default}, \texttt{palatino}, \texttt{charter}, \texttt{sans} & \texttt{default} \\
\texttt{showsolutions} & \texttt{true}, \texttt{false} & \texttt{true} \\
\texttt{parskip} & \texttt{true}, \texttt{false} & \texttt{true} \\
\bottomrule
\end{tabular}
\caption{Class options and parameter values}
\end{table}

\section{Course Metadata Commands}

Specify your course metadata in the document preamble:
\begin{codebox}[Metadata Setup]
\begin{lstlisting}[language=TeX]
\course{Data Structures & Algorithms}
\coursecode{CS 201}
\professor{Dr. Alex Morgan}
\institution{Department of Computer Science}
\term{Fall 2026}
\docsubtitle{Comprehensive Lecture Notes}
\title{CS 201: Class Notes}
\end{lstlisting}
\end{codebox}

\section{Lecture Tracking System}

Use the \texttt{\textbackslash lecture} command to divide course notes into distinct lecture sessions:
\begin{codebox}[Lecture Macro]
\begin{lstlisting}[language=TeX]
\lecture[September 2, 2026]{1}{Introduction to Algorithm Analysis}
\end{lstlisting}
\end{codebox}

\section{Pedagogical Environments}

\begin{definition}[Mathematical Callout Environments]
\texttt{modernclassnotes} provides structured environments for mathematical and academic content:
\begin{itemize}
    \item \texttt{theorem}, \texttt{lemma}, \texttt{proposition}, \texttt{corollary}, \texttt{proof}
    \item \texttt{definition}, \texttt{example}, \texttt{remark}, \texttt{warning}, \texttt{tip}, \texttt{takeaways}
    \item \texttt{exercise}, \texttt{solution}, \texttt{codebox}
\end{itemize}
\end{definition}

\begin{example}[Sample Theorem Environment]
\begin{lstlisting}[language=TeX]
\begin{theorem}[Euler's Identity]
For any real number $\theta$, $e^{i\theta} = \cos\theta + i\sin\theta$.
\end{theorem}
\end{lstlisting}
\end{example}

\section{Open-Source \& CTAN Distribution Guidelines}

To distribute this package under an open-source initiative:
\begin{enumerate}
    \item Maintain the LPPL 1.3c / MIT dual licensing terms in \texttt{LICENSE}.
    \item Provide compiled sample PDFs in \texttt{examples/} for GitHub previews.
    \item Execute \texttt{make ctan} to bundle a clean \texttt{.zip} release file ready for CTAN submission.
\end{enumerate}

\end{document}
